\documentclass[11pt]{article}
\usepackage[margin=1in]{geometry}
\usepackage{amsmath,amssymb}
\usepackage[T1]{fontenc}
\usepackage{hyperref}

\title{Literature status note for arXiv:1209.0055\\
Two-dimensional Mazur--Ulam-property subcase}
\author{agent\_lane\_16}
\date{2026-06-29}

\begin{document}
\maketitle

\section*{Status}

\textbf{literature\_implied\_answer (partial subcase).}

\section*{Source problem}

The source paper is Dongni Tan and Rui Liu, \emph{A Note on The
Mazur-Ulam Property of Almost-CL-spaces}, arXiv:1209.0055.  In the
introduction they recall Tingley's isometric extension problem as
Problem~1:

\begin{quote}
Let $X$ and $Y$ be normed spaces with unit spheres $S_X$ and $S_Y$.
Assume that $T:S_X\to S_Y$ is a surjective isometry.  Does there exist
a linear isometry $\widetilde T:X\to Y$ such that
$\widetilde T|_{S_X}=T$?
\end{quote}

They then formulate the Mazur--Ulam property and state that the problem
is still open in general, even in two dimensions.  The same paper proves
positive results for spaces with the $(T)$-property, almost-CL spaces
admitting smooth points, one hexagonal two-dimensional norm, and some
stability constructions; those same-paper results are not the basis of
this status note.

\section*{Supporting theorem}

The supporting paper is Taras Banakh, \emph{Every 2-dimensional Banach
space has the Mazur--Ulam property}, arXiv:2103.09268; Linear Algebra
Appl. 632 (2022), 268--280.

Its abstract states that every isometry between the unit spheres of
2-dimensional Banach spaces extends to a linear isometry and that this
resolves Tingley's problem in the class of 2-dimensional Banach spaces.
The main theorem, Theorem~1, states:
\[
  \text{Every 2-dimensional Banach space has the Mazur--Ulam property.}
\]
This theorem answers the two-dimensional subcase that Tan--Liu described
as still open in 2012.

\section*{Scope}

This packet does not claim a full solution of Tingley's problem or of the
general Cheng--Dong Mazur--Ulam question asking whether every Banach
space has the Mazur--Ulam property.  The supporting theorem settles only
the 2-dimensional case.  Later literature still describes the problem as
open in finite dimensions at least three, while citing Banakh for the
2-dimensional positive answer.

The relation is an agent-identified implication.  Banakh's paper cites
Cheng--Dong for the Mazur--Ulam-property formulation and explicitly
solves the named Tingley problem in dimension two, but it does not appear
to cite Tan--Liu's almost-CL paper directly.

\section*{Search evidence}

The run indexes were searched for arXiv:1209.0055 and the keywords
``Mazur-Ulam property'', ``Tingley'', ``almost-CL'', ``MUP'', and
``sphere isometry''.  No exact duplicate packet for arXiv:1209.0055 was
found, although nearby Tingley/MUP packets already exist for other
subclasses.  External web search on 2026-06-29 and local parsed-source
inspection identified arXiv:2103.09268 as the decisive answer for the
two-dimensional subcase.

\end{document}
